\documentclass[handout]{beamer}
\mode<presentation>
\usepackage{xcolor}
\usepackage[toc]{multitoc}
\usepackage{color}
\usepackage{graphicx}
\graphicspath{{../buku_ajar/buku_ajar_logika_matematika/figures/}}
\usepackage{tikz}
\usetikzlibrary{shapes,arrows,positioning,calc}
\usepackage{listings}
\usepackage{lmodern}
\usepackage{amsmath}
\usetheme[secheader]{Boadilla}
\usecolortheme{whale}
\setbeamertemplate{caption}[numbered]
\setbeamertemplate{bibliography item}[text]
\setbeamertemplate{navigation symbols}{}

\newcommand{\logic}[1]{\textbf{\textit{#1}}}

\theoremstyle{definition}
\newtheorem{definisi}[theorem]{Definisi}
\theoremstyle{example}
\newtheorem{contoh}[theorem]{Contoh}

\renewcommand*{\multicolumntoc}{2}

\setbeamertemplate{footline}
{
  \leavevmode%
  \hbox{%
  \begin{beamercolorbox}[wd=.333333\paperwidth,ht=2.25ex,dp=1ex,center]{author in head/foot}%
    \usebeamerfont{author in head/foot}\insertshortauthor%
  \end{beamercolorbox}%
  \begin{beamercolorbox}[wd=.433333\paperwidth,ht=2.25ex,dp=1ex,center]{title in head/foot}%
    \usebeamerfont{title in head/foot}\insertshorttitle
  \end{beamercolorbox}%
  \begin{beamercolorbox}[wd=.233333\paperwidth,ht=2.25ex,dp=1ex,right]{date in head/foot}%
    \usebeamerfont{date in head/foot}\insertshortdate{}\hspace*{2em} \insertframenumber{} / \inserttotalframenumber\hspace*{2ex} 
  \end{beamercolorbox}}%
  \vskip0pt%
}

\subtitle[Logika Matematika]{Logika Matematika}
\title[Validasi Algoritma]{Argumen Formal dan Validasi Algoritma}
\author{Cecep Suwanda, S.Si., M.Kom.}
\date{Maret 2025}
\institute{Universitas Bale Bandung}

\begin{document}

\maketitle

\section*{Daftar Isi}
\begin{frame}
\frametitle{Daftar Isi}
\tableofcontents
\end{frame}

% ============================================================
% SECTION 14-1: Pre/Post
% ============================================================
\section{Prekondisi dan Postkondisi}

\begin{frame}
\frametitle{Prekondisi dan Postkondisi}
\begin{definisi}[Prekondisi]
Predikat yang harus dipenuhi \textbf{sebelum} blok kode dijalankan.
\end{definisi}
\begin{definisi}[Postkondisi]
Predikat yang harus terpenuhi \textbf{setelah} blok selesai (jika pre terpenuhi).
\end{definisi}
``Kontrak'' antara pemanggil dan implementasi.
\end{frame}

\begin{frame}
\frametitle{Contoh Spesifikasi}
\begin{itemize}
  \item \texttt{sqrt}($x$): pre $x \geq 0$; post $r \geq 0$ dan $r^2 = x$
  \item \texttt{divide}($a,b$): pre $b \neq 0$; post $a = bq + r$, $0 \leq r < |b|$
\end{itemize}
Melanggar pre $\Rightarrow$ perilaku tidak dijamin.
\end{frame}

% ============================================================
% SECTION 14-2: Hoare
% ============================================================
\section{Hoare Triple}

\begin{frame}
\frametitle{Hoare Triple}
\begin{definisi}[Hoare Triple]
$\{p\}\ S\ \{q\}$ valid jika: setiap eksekusi $S$ dari keadaan yang memenuhi $p$ yang \textbf{berterminasi} menghasilkan keadaan yang memenuhi $q$.
\end{definisi}
Logika Hoare (Tony Hoare, 1969).
\end{frame}

\begin{frame}
\frametitle{Aturan Assignment}
$\{q[E/x]\}\ x := E\ \{q\}$\\
Contoh: $\{x = 5\}\ x := x+1\ \{x = 6\}$ (substitusi mundur).
\end{frame}

\begin{frame}
\frametitle{Komposisi dan If}
\textbf{Komposisi:} $\dfrac{\{p\}\ S_1\ \{r\} \quad \{r\}\ S_2\ \{q\}}{\{p\}\ S_1; S_2\ \{q\}}$\\
\textbf{If:} kedua cabang dari $p \wedge B$ dan $p \wedge \neg B$ ke post $q$ yang sama.
\end{frame}

% ============================================================
% SECTION 14-3: Loop invariant
% ============================================================
\section{Loop Invariant}

\begin{frame}
\frametitle{Tiga Syarat Loop Invariant}
\begin{definisi}[Loop Invariant]
Predikat $I$ untuk loop \texttt{while $B$ do $S$}:
\begin{enumerate}
  \item \textbf{Inisialisasi}: $I$ benar sebelum iterasi pertama
  \item \textbf{Pemeliharaan}: $\{I \wedge B\}\ S\ \{I\}$
  \item \textbf{Terminasi}: $I \wedge \neg B$ mengimplikasikan postkondisi
\end{enumerate}
\end{definisi}
\end{frame}

\begin{frame}
\frametitle{Aturan While}
$\dfrac{\{I \wedge B\}\ S\ \{I\}}{\{I\}\ \texttt{while } B \texttt{ do } S\ \{I \wedge \neg B\}}$
\end{frame}

\begin{frame}
\frametitle{Analogi dengan Induksi}
\begin{center}
\small
\begin{tabular}{|l|l|}
\hline
\textbf{Induksi} & \textbf{Loop invariant} \\ \hline
Basis $P(n_0)$ & Inisialisasi $I$ \\ \hline
$P(k) \Rightarrow P(k+1)$ & Pemeliharaan \\ \hline
Kesimpulan $P(n)$ & $I \wedge \neg B$ \\ \hline
\end{tabular}
\end{center}
\end{frame}

\begin{frame}
\frametitle{Contoh: Penjumlahan Array}
\texttt{sum := 0; i := 0; while i < n do sum := sum + A[i]; i := i+1}\\
$I$: $sum = \sum_{j=0}^{i-1} A[j]$ dan $0 \leq i \leq n$.\\
Setelah loop: $i = n \Rightarrow$ post $sum = \sum_{j=0}^{n-1} A[j]$.
\end{frame}

% ============================================================
% SECTION 14-4: Parsial vs total
% ============================================================
\section{Kebenaran Parsial dan Total}

\begin{frame}
\frametitle{Kebenaran Parsial}
$\{p\}\ S\ \{q\}$: jika $p$ dan $S$ \textbf{berterminasi}, maka $q$.\\
Tidak menjamin terminasi (infinite loop: vakum terpenuhi).
\end{frame}

\begin{frame}
\frametitle{Kebenaran Total}
$[p]\ S\ [q]$: dari $p$, $S$ \textbf{pasti} berterminasi dan $q$ terpenuhi.\\
Total = parsial + bukti terminasi.
\end{frame}

\begin{frame}
\frametitle{Fungsi Pengurangan (Variant)}
\begin{definisi}[Fungsi pengurangan]
$V \geq 0$ selama $B$ benar; $V$ turun ketat tiap iterasi $\Rightarrow$ loop berhenti.
\end{definisi}
Contoh penjumlahan array: $V = n - i$.
\end{frame}

\begin{frame}
\frametitle{Ringkasan}
\begin{center}
\scriptsize
\begin{tabular}{|l|l|l|}
\hline
 & \textbf{Parsial} $\{p\}S\{q\}$ & \textbf{Total} $[p]S[q]$ \\ \hline
Terminasi & Tidak dijamin & Dijamin \\ \hline
Post & Jika berhenti & Selalu \\ \hline
Bukti & Invariant & Invariant + variant \\ \hline
\end{tabular}
\end{center}
\end{frame}

% ============================================================
% SECTION 14-5: Linear search
% ============================================================
\section{Pencarian Linear}

\begin{frame}
\frametitle{Loop Invariant Pencarian Linear}
$I$: untuk semua $j$, $0 \leq j < i \Rightarrow A[j] \neq key$, dan $0 \leq i \leq n$.\\
Inisialisasi $i=0$: vakum.\\
Terminasi normal: $i=n \Rightarrow$ key tidak ada, return $-1$.
\end{frame}

\begin{frame}
\frametitle{Terminasi}
$V = n - i$; tiap iterasi (tanpa return) $i$ naik $\Rightarrow$ terminasi.
\end{frame}

% ============================================================
% SECTION 14-6: Insertion sort
% ============================================================
\section{Insertion Sort}

\begin{frame}
\frametitle{Spesifikasi Pengurutan}
Post: (1) $A[0] \leq \cdots \leq A[n-1]$ terurut;\\
(2) $A$ = \textbf{permutasi} array masukan (penting!).
\end{frame}

\begin{frame}
\frametitle{Invariant Loop Luar}
Pada awal iterasi $j$: $A[0..j-1]$ terurut dan permutasi elemen asli $A[0..j-1]$.\\
$j=n$: seluruh array terurut dan permutasi.
\end{frame}

\begin{frame}
\frametitle{Loop Dalam}
Invariant: $A[i+2..j]$ berisi elemen $>$ key terurut; key akan ke $A[i+1]$.\\
Variant: $V = i + 1$ (turun saat $i$ turun).
\end{frame}

% ============================================================
% SECTION 14-7: Industri
% ============================================================
\section{Verifikasi Formal}

\begin{frame}
\frametitle{Testing vs Verifikasi Formal}
\begin{itemize}
  \item Testing: sampel kasus; menunjukkan \textbf{keberadaan} bug
  \item Formal: bukti untuk semua input; ketiadaan bug (dalam model)
\end{itemize}
Dijkstra: testing tidak pernah membuktikan ketiadaan bug.
\end{frame}

\begin{frame}
\frametitle{Mission-Critical}
Aviasi, medis, antariksa (Ariane 5), radioterapi (Therac-25), kriptografi.
\end{frame}

\begin{frame}
\frametitle{Alat Bantu}
Static analyzers; model checkers (SPIN, NuSMV); theorem provers (Coq, Isabelle, Lean); Design by Contract (Eiffel, JML).
\end{frame}

\begin{frame}
\frametitle{Konvergensi Materi Buku}
Proposisi/predikat $\to$ spesifikasi; pembuktian $\to$ parsial; induksi $\to$ invariant.
\end{frame}

% ============================================================
% RANGKUMAN
% ============================================================
\section{Rangkuman}

\begin{frame}
\frametitle{Rangkuman}
\begin{itemize}
  \item Pre/post = kontrak; Hoare $\{p\}S\{q\}$
  \item Invariant: inisialisasi, pemeliharaan, terminasi
  \item Parsial vs total; variant untuk terminasi
  \item Pengurutan: terurut + permutasi
\end{itemize}
\end{frame}

\begin{frame}
\frametitle{Kata Kunci}
\begin{center}
\small
prekondisi $\bullet$ postkondisi $\bullet$ Hoare Triple $\bullet$ assignment $\bullet$ loop invariant $\bullet$ kebenaran parsial $\bullet$ kebenaran total $\bullet$ fungsi pengurangan $\bullet$ verifikasi formal
\end{center}
\end{frame}

\end{document}
